\section{Rによる群間計画の分析}
\label{b19_BetweenDesign_onR}

\subsection{授業内容}
\subsubsection{概要}
前回学んだ群間計画(Between-subjects Design)の理論を踏まえ，本コマでは統計環境Rを用いた実践的な分析方法について学ぶ。一要因および二要因の群間計画データを実際に分析し，分散分析表の読み方や結果の解釈方法を習得する。基本的な分析には\verb|aov()|関数と\verb|lm()|関数を用い，それらの出力結果からSS(平方和)，df(自由度)，MS(平均平方)，F値などの情報をどのように読み取るかを学ぶ。また，分析結果を適切に可視化する方法についても実習を通じて理解する。

\subsubsection{コマ主題細目}
\begin{description}
	\item[データの準備と基本操作] 分散分析を行うためのデータ形式について理解し，適切なデータの読み込み方法と前処理を学ぶ。特にRでの分散分析では，カテゴリカル変数を因子型(factor)として扱うことが重要である。csvファイルやExcelファイルからのデータの読み込み方法，データフレームの構造確認，因子型への変換方法，データの要約統計量の算出方法などの基本操作を復習する。また，欠損値の処理方法についても学び，分析前の適切なデータチェックの重要性を理解する。
	      \begin{flushright}
		      $\to$ \textcite{kosugi2019}のP.46--70，\textcite{Kinosady2021}のP.28--42.
	      \end{flushright}

	\item[一要因群間計画の分析] Rの\verb|aov()|関数と\verb|lm()|関数を用いた一要因分散分析の実行方法を学ぶ。公式表記(formula notation)の書き方(例：\verb|y ~ group|)，関数の使い方，出力結果の読み方について理解する。\verb|summary()|関数を用いた分散分析表(ANOVA table)の表示方法と，そこに含まれるSS(平方和)，df(自由度)，MS(平均平方)，F値，p値の意味を確認する。また，\verb|TukeyHSD()|関数を用いた多重比較の方法や，効果量($\eta^2$，$\omega^2$など)の算出方法についても学ぶ。実際のデータセットを用いた演習を通じて，一要因群間計画の分析手順と結果の解釈を身につける。
	      \begin{flushright}
		      $\to$ \textcite{kosugi2019}のP.146--152，\textcite{Hashimoto201607}のP.26--55.
	      \end{flushright}

	\item[二要因群間計画の分析] 二要因分散分析の実行方法と結果の解釈について学ぶ。\verb|aov()|関数における二要因の公式表記(例：\verb|y ~ factor1 * factor2|)の意味と使い方，交互作用項の指定方法を理解する。分散分析表から主効果と交互作用の検定結果を読み取る方法，有意な交互作用がある場合の解釈と追加分析(単純主効果の検定など)の必要性について学ぶ。また，二要因分散分析の結果を適切に報告する方法についても理解する。実際のデータセットを用いた演習を通じて，二要因群間計画の分析手順と結果の解釈を身につける。
	      \begin{flushright}
		      $\to$ \textcite{kosugi2019}のP.153--160，\textcite{Hashimoto201607}のP.57--90.
	      \end{flushright}

	\item[分析結果の可視化] 分散分析の結果を適切に可視化する方法について学ぶ。Rの基本グラフィックス機能やggplot2パッケージを用いた群間の比較を示す棒グラフやボックスプロット，交互作用を示す交互作用プロットの作成方法を理解する。エラーバーの追加方法，軸ラベルやタイトルの設定，凡例の調整など，読みやすく情報量の多いグラフを作成するためのテクニックを学ぶ。また，可視化することの重要性と，異なるグラフが伝える情報の違いについても理解する。実際のデータセットを用いた演習を通じて，分析結果の効果的な可視化方法を身につける。
	      \begin{flushright}
		      $\to$ \textcite{Kinosady2021}のP.126--138，\textcite{kosugi2019}のP.90--105.
	      \end{flushright}

\end{description}
\subsubsection{キーワード}
\begin{itemize}
	\item \verb|aov()|関数と\verb|lm()|関数
	\item 分散分析表の読み方(SS, df, MS, F値)
	\item 多重比較と下位検定
	\item 交互作用の可視化
	\item ggplot2による結果の表現
\end{itemize}
\subsection{授業情報}
\paragraph{コマの展開方法}
演習
\paragraph{標準シラバスにおける位置づけ}
科目番号5 心理学統計法；2統計に関する基礎的な知識；C/D/E エクセル，R，SPSS入門
\subsubsection{予習・復習課題}
\paragraph{予習}
前回学んだ群間計画(Between Design)の基本概念と分散分析の理論を復習しておくこと。特に，平方和の分解，F値の意味，主効果と交互作用の概念について理解を確認しておくとよい。また，RとRStudioの基本操作，特にデータの読み込み方法やデータフレームの操作方法について再確認しておくこと。必要に応じて，第3回，第4回の授業で学んだR操作についての資料を見直しておくとよい。

\paragraph{復習}
授業で学んだ内容を定着させるために，以下の課題に取り組むこと：

\begin{enumerate}
\item 提供されたデータセットを用いて，一要因群間計画の分散分析を実行し，結果を適切に解釈する
\item 提供されたデータセットを用いて，二要因群間計画の分散分析を実行し，主効果と交互作用の結果を解釈する
\item 分散分析の結果を適切に可視化し，群間の差や交互作用を示すグラフを作成する
\item \verb|TukeyHSD()|関数を用いて多重比較を行い，どの群間に有意差があるかを特定する
\item 上記の分析結果をRmarkdownを用いてレポート形式にまとめる
\end{enumerate}

これらの課題をRmarkdownファイルにまとめ，期日までに提出すること。時間内に終わらなかった場合は，次回授業までに完成させておくこと。分析結果の適切な解釈と可視化を重視し，統計的に有意な結果だけでなく，効果量や信頼区間についても言及するよう心がけること。不明点があれば，\textcite{kosugi2019}や\textcite{Hashimoto201607}を参照するか，オフィスアワーを利用して質問すること。